\section{Introducción}

El diseño e implementación de un sistema autónomo para el control de acceso e
identificación vehicular en parqueaderos urbanos requiere una fundamentación 
teórica multidisciplinaria dentro de la visión por computadora y el aprendizaje
profundo. Este capítulo establece las bases conceptuales, los modelos
matemáticos y las evoluciones algorítmicas que sustentan el flujo de
procesamiento completo del proyecto: desde la captura del fotograma digital en
las barreras de ingreso hasta la determinación automatizada de los eventos de
acceso.

El desarrollo temático de este marco teórico sigue una estructura jerárquica
descendente (\textit{top-down}) organizada en seis secciones que avanzan desde
los fundamentos físicos de la percepción visual hasta la evaluación cuantitativa
del sistema.

La primera sección examina los principios de la \textbf{visión artificial clásica
y el procesamiento de imágenes}, abarcando las representaciones matriciales, los
espacios cromáticos, el filtrado espacial y las transformaciones morfológicas
orientadas al preprocesamiento de matrículas, así como sus limitaciones físicas
ante variaciones de iluminación, sombras y oclusiones en entornos no
controlados. La segunda establece los \textbf{fundamentos del aprendizaje
automático y profundo}, desde el paradigma supervisado y la partición de datos
hasta las redes neuronales convolucionales, el aprendizaje por transferencia y
las técnicas de aumento de datos.

Sobre esa base, la tercera sección aborda la \textbf{detección de objetos},
justificando la adopción de arquitecturas de una sola etapa ---como la familia
YOLO--- frente a las alternativas de dos etapas, y analizando los cuellos de
botella computacionales que condicionan la inferencia en tiempo real. La cuarta
se concentra en la \textbf{percepción vehicular y el reconocimiento de
matrículas}, cubriendo la clasificación de atributos, la localización de objetos
pequeños y los modelos de secuencia con alineación temporal. La quinta desarrolla
el \textbf{seguimiento vehicular y la determinación de eventos}, desde la
formulación del seguimiento multiobjeto hasta la interpretación de trayectorias.
Finalmente, la sexta sección cierra con la \textbf{evaluación y optimización de
los modelos}, definiendo las métricas de desempeño y las técnicas de compresión
que habilitan su despliegue en hardware de borde.
